\section{Feedback Edge Set}

A \emph{feedback edge set} of a graph is a set of edges that when
removed renders the graph acyclic. It is well-known that if $G$ is a connected
undirected graph on $n$ vertices and $m$ edges, then all minimal feedback edge
sets of $G$ have size $k=m-n+1$. Indeed, such a set can be constructed in
polynomial time by repeatedly locating a cycle in the graph and selecting an
arbitrary edge of the cycle to insert into the feedback edge set.


In this section we will consider \SF parameterized by the feedback edge set of
the input graph, which we will denote by $k$. Unlike the vertex cover section,
here our main result is positive: we show that \SF can be solved optimally in
time $2^{O(k)}n^{O(1)}$, that is, in time single-exponential in the parameter.
Since we are able to achieve a single-exponential dependence, it is
straightforward to see that this is optimal under the ETH.

\begin{theorem}\label{thm:feseth} %
If there is an algorithm solving \ST in time
$2^{o(k)}n^{O(1)}$, where $k$ is the feedback edge set of the input, then the
ETH is false. \end{theorem}

%

\begin{proof}%
The proof follows from the sparsification lemma of Impagliazzo,
Paturi, and Zane \cite{ImpagliazzoPZ01} composed with the standard reduction
proving that \ST is NP-complete. We sketch the details. Suppose we are given a
\textsc{3-SAT} formula $\phi$ with $n$ variables and $m$ clauses. The
sparsification lemma shows that in order to disprove the ETH it is sufficient
to show that we can decide if $\phi$ is satisfiable in time $2^{o(n)}$ under
the restriction that $m=\Theta(n)$. We edit $\phi$ to obtain an equisatisfiable
formula $\phi'$ where every variable appears at most $3$ times (for each
variable $x$ appearing $f>3$ times, we replace each occurrence of $x$ with $f$
fresh variables $x_1,\ldots, x_f$ and add the clauses $(x_1\to x_2)\land
(x_2\to x_3)\land\ldots (x_f\to x_1)$). The new formula $\phi'$ has $n'=O(n)$
variables and $m'=O(n)$ clauses. We now execute the chain of reductions showing
that \ST is NP-hard (e.g. from  \cite{karp1975computational}), which produce an
instance on a graph $G=(V,E)$ with $|E|=O(m')$, therefore, $|E|=O(n)$. The new
instance has feedback edge set size $k<|E|$, therefore an algorithm solving the
new instance in time $2^{o(k)}|V|^{O(1)}$ would falsify the ETH.  \end{proof}
    
%

\begin{comment}
Before we proceed, let us give a high-level overview of the algorithm. We will
begin by applying a simplification procedure which eliminates degree one
vertices from the graph (\cref{lem:leaf}). Via simple counting arguments we
then observe that the resulting instance contains only $O(k)$ \emph{special}
vertices, that is vertices that have degree at least $3$ or are incident on an
edge of the feedback edge set (\cref{lem:deg3}).  Furthermore, these vertices
are connected via $O(k)$ \emph{topo-edges}, that is, paths that go through
non-special degree two vertices (\cref{lem:topo}). We then guess which
topo-edges are fully included in the solution (this gives the $2^{O(k)}$ factor
in the running time). Once we have this information, we can infer how the
special vertices are partitioned into connected components in the optimal
solution. 

What remains is to decide what to do at the interior of topo-edges which are
not fully used. For this we do a case analysis: for a topo-edge
whose endpoints are in the same component of the optimal, things are intuitively
simple because a demand that needs to be routed from inside the topo-edge to
the outside can pass through either endpoint and this allows us to isolate the
demands inside the topo-edge and solve locally for optimality; for a topo-edge
whose endpoints are in distinct components, all demands that are internal to
the topo-edge have a unique feasible way to be satisfied, so we select this; we
therefore arrive at a situation where all remaining demands are between
internal vertices of distinct topo-edges. In this setting, the only scenario
where we may have more than one reasonable way to satisfy a demand is if we
have two topo-edges both of which have one endpoint in a component $C_1$ of the
optimal and the other in a component $C_2$ of the optimal. For each pair of
components $C_1,C_2$ that contain special vertices, we therefore consider all
topo-edges with one endpoint in $C_1$ and the other in $C_2$. Intuitively, this
reduces to solving \SF on a very special graph class: we have two vertices and
a large number of parallel paths connecting them, with all demands having
endpoints in distinct paths. Despite its apparent simplicity, this special case
is still non-trivial, but we manage to solve such instances in polynomial time
via a reduction to \textsc{Min Cut}, thus completing the algorithm.

\end{comment}

Let us now proceed to the detailed presentation of the algorithm. Suppose that
we are given 
%
%
a budget~$b$ and we want to decide if there exists a \SF solution $F$ such
that $\cost(F)\le b$.
%
We start by applying a simple reduction rule.

\begin{description} 
\item[Rule 3:]\customlabel{rule1fes}{Rule 3} Suppose we have a \SF instance on
graph $G$ with weight function $w$ and budget~$b$, such that a vertex $u\in V$
has degree $1$.  If $u\not\in R$, then delete $u$.  If $u\in R$, let $v$ be the
unique neighbor of $u$. Then  set $b':=b-w(uv)$, delete $u$ from the graph and
the demand $\{u,v\}$ from $D$ if it exists, and replace, for each $x\in
V\setminus\{u,v\}$ such that $\{u,x\}\in D$ the demand $\{u,x\}$ with the
demand $\{v,x\}$.
\end{description}

\begin{lemma}\label{lem:leaf} %
\ref{rule1fes} is safe. 
\end{lemma}

%
\begin{proof}%
If $u\not\in R$ then no optimal solution contains the edge $uv$,
so it is safe to delete $u$. If $u\in R$ then all feasible solutions contain
the edge $uv$. \end{proof}
%

Observe that if we apply \ref{rule1fes} exhaustively, then the minimum degree of the graph is~$2$. As we show next,
relatively few vertices can have higher degree.

\begin{lemma}\label{lem:deg3} %
Suppose we have a \SF instance with feedback edge
set of size $k$ and minimum degree at least $2$.  Then $G$ contains at most
$2k$ vertices of degree at least $3$. 
\end{lemma}

%
\begin{proof}%
We observe that if our graph has a feedback edge set of size
$k$, then $m= k+n-c$, where $c$ is the number of connected components of $G$.
This implies that $\sum_{v\in V}d(v) = 2m = 2k+2n-2c$. Let $V_2$ be the set of
vertices of degree exactly $2$ and $V_3=V\setminus V_2$ be the set of vertices
of degree at least $3$. We have $\sum_{v\in V}d(v) \ge 2|V_2|+3|V_3| = 2n
+|V_3|$. We conclude that $|V_3|\le 2k-2c$.  \end{proof}
%


In the remainder we will assume that we have a \SF instance $G=(V,E)$ with a
feedback edge set $H\subseteq E$ of size $k$, to which \ref{rule1fes} can no longer be
applied. We will say that a vertex $v$ is \emph{special} if $v$ is incident on
an edge of $H$ or $v$ has degree at least $3$. By \cref{lem:deg3} we know that
$G$ contains at most $4k$ special vertices.

We define a \emph{topological edge} (topo-edge for short) as follows: a path
$P$ in $G$ is a topological edge if the two endpoints of $P$ are special
vertices and all internal vertices of $P$ are non-special. Note that by this
definition, all edges of $H$ form topo-edges, since the endpoints of such edges
are special.  We observe the following:

\begin{lemma}\label{lem:topo} %
Suppose we have a graph $G$ with
feedback edge set of size $k$ and minimum degree at least~$2$. Then $G$
contains at most $5k$ topological edges.  
\end{lemma}

%

\begin{proof}%
Let $V_s$ be the set of special
vertices and $V_t=V\setminus V_s$. If we have more than $5k$ topological edges
in $G$ then $\sum_{v\in V_s}d(v)\ge 10k$. This is because each topological edge
contributes at least $2$ in the sum $\sum_{v\in V_s}d(v)$.  On the other hand,
if $c$ is the number of connected components of $G$, we have $2n+2k-2c = 2m =
\sum_{v\in V}d(v) = \sum_{v\in V_s}d(v) + \sum_{v\in V_t} d(v) = \sum_{v\in
V_s} d(v) +2|V_t|$.  However, $|V_t|\ge n-4k$ by \cref{lem:deg3} and the fact
that at most $2k$ vertices are incident on $H$.  Hence, $2n+2k-2c\ge \sum_{v\in
V_s}d(v) +2n-8k$.  This implies that $\sum_{v\in V_s}d(v)<10k$.  Hence, it is
impossible to have more than $5k$ topological edges. \end{proof}
%

We are now ready to state the main algorithmic result of this section.

\begin{theorem} 
There is an algorithm that solves \SF on instances with $n$
vertices and a feedback edge set of size $k$ in $2^{O(k)}n^{O(1)}$ time. 
\end{theorem}

\begin{proof}
%
%
Call the set of special vertices $V_s$ and let $V_t=V\setminus V_s$. For the rest of
this proof and for the sake of the analysis, fix an optimal solution $F^\star$.

To begin, we guess which of the $5k$ topological edges according to \cref{lem:topo} are fully used in the
optimal solution. To be more precise, we will say that a topo-edge $P$ is fully
used in $F^\star$ if all edges of the path $P$ are contained in~$F^\star$. This gives
$2^{5k}$ possibilities. In the remainder we will assume that we have correctly
guessed the set of topo-edges which are fully used in $F^\star$.

We now observe that for any two vertices $u,v\in V_s$ we have enough
information to deduce whether~$u,v$ are in the same connected component of
$F^\star$. More precisely, we construct an auxiliary graph $G_s$ with vertex
set $V_s$ that contains an edge between two vertices $u,v\in V_s$ if there
exists a fully used topo-edge whose endpoints are $u,v$.  We now claim that two
vertices $u,v\in V_s$ are in the same component of $F^\star$ if and only if
$u,v$ are in the same connected component of $G_s$. Indeed, if two vertices
$u,v$ are in the same component of $G_s$, then clearly there is a path
connecting them in $F^\star$ going through fully used topo-edges; conversely if
$u,v$ are in the same component of $F^\star$ and the path connecting them goes
through the special vertices $u_1=u, u_2,\ldots, u_\ell=v$ (and all other
vertices are non-special), then the path $u_1,\ldots, u_\ell$ also exists in
$G_s$, as the topo-edge connecting $u_i, u_{i+1}$ must be fully used.

Because of the above, we can now assume that we have a partition $\rho$ of $V_s$
such that $u,v$ are in the same set of $\rho$ if and only if $u,v$ are in the
same connected component of $F^\star$. Notice that this implies that we can remove
from the instance all demands $\{u,v\}\in D$ such that $u,v\in V_s$: if $u,v$ are in
the same set of $\rho$ the demand is automatically satisfied by our guess of the
fully used topo-edges; while if $u,v$ are in distinct sets of $\rho$, we know
that our guess is incorrect and we reject the current instance. Every remaining
demand of our instance is therefore incident on at least one non-special
vertex.


What remains is to decide for topo-edges which are not fully used, which of
their incident edges belong in $F^\star$. Note that this is trivial for topo-edges
consisting only of a single edge, since fully using such a topo-edge is
equivalent to placing the corresponding edge in the solution. We therefore
focus on topo-edges which contain at least one internal (non-special) vertex. 

For this we proceed in several steps.  First, suppose we have a non-fully-used
topo-edge $P$ whose endpoints are adjacent to $u,v\in V_s$ such that $u,v$ are
in the same component of $F^\star$. We edit the instance so that demands with one
endpoint in the interior of $P$ also have their other endpoint in $P$. More
precisely, for each demand $\{x,y\}\in D$ such that $x$ is an internal vertex
of $P$ and $y\not\in P$, we remove $\{x,y\}$ from $D$ and replace it with the
demands $\{x,u\}$ and $\{y,u\}$. It is not hard to see that this is safe,
because any path satisfying the demand $\{x,y\}$ would have to go either
through $u$ or through $v$, but $u,v$ are in the same component of $F^\star$ thanks
to other, fully-used topo-edges, so routing the demand through $u$ or $v$ is
the same.

Consider then a topo-edge $P$ whose endpoints $u,v$ are in the same component
of $F^\star$ and where all demands with one endpoint in an internal vertex of $P$
have the other endpoint in $P$. We simplify the instance by branching: select
an edge $e\in P$, delete $e$ from the instance, and apply \ref{rule1fes} exhaustively
on internal vertices of $P$, until all such vertices are removed.  Since we
have guessed that $P$ is not fully used, at least one of the instances we
produced is equivalent to the original, that is, at least one choice of edge to
delete indeed deletes an edge not used by the optimal solution. It may seem
that since we are branching on $n$ possibilities, this branching will lead to a
running time of $n^k$.  However, we observe that after removing any edge of $P$
and exhaustively applying \ref{rule1fes} we obtain instances which have (i)
the same graph, as all internal vertices of $P$ have been deleted and all other
vertices are unchanged (ii) the same set of demands, as all demands with one
endpoint in an internal vertex of $P$ have either been removed or replaced with
the demand $\{u,v\}$ (which is satisfied by the fully used topo-edges, so can
be removed) and other demands are unchanged.  Hence, among the at most $n$
instances this branching produces, it suffices to select the one with the
maximum remaining budget and solve that, to decide if the original is a Yes
instance. In other words, the branching procedure of this paragraph is a
polynomial-time reduction rule which allows us to eliminate all topo-edges
whose endpoints are in the same component of~$F^\star$.

In the remainder we thus assume that every topo-edge $P$ that is not fully used
has endpoints $u,v\in V_s$ which are in distinct components of $F^\star$. Next, we
deal with the case of ``internal'' demands. Suppose that there exists a
topo-edge $P$ with endpoints $u,v\in V_s$ that contains an internal demand,
that is, there exist $w_1,w_2\in P\setminus\{u,v\}$ such that $\{w_1,w_2\}\in
D$.  Then, all edges in the path from $w_1$ to $w_2$ in $P$ must belong in
$F^\star$, because every other solution that connects $w_1$ to $w_2$ would put
$u,v$ in the same component. We can therefore contract all the edges of the
path from $w_1$ to $w_2$ and adjust our budget and our demands
accordingly: we decrease our budget by the total cost of the edges of
the path from $w_1$ to~$w_2$, we remove all demands that have both endpoints in
that path, and for demands that have one endpoint in that path, we replace that
endpoint by the vertex that results from the contraction of the path.

We now arrive at the case where the endpoints of each topo-edge are adjacent to
vertices from distinct components of $F^\star$ and demands with one endpoint in the
interior of a topo-edge have the other endpoint outside of the topo-edge or in
$V_s$.

We distinguish several cases:

\begin{enumerate}

\item There exists a topo-edge $P$ adjacent to $u,v\in V_s$, an internal vertex
$w\in P\setminus\{u,v\}$ and a vertex $w'\in V_s$ such that $\{w,w'\}\in D$. If
$w'$ is in the same component of $F^\star$ as $u$ (respectively $v$), we include in
the solution all edges in the path in $P$ from $w$ to $u$ (respectively $v$),
contract the selected edges and update our budget and demands accordingly, as above. If
$w',u,v$ are in distinct components of~$F^\star$, then we conclude that the current
guess is incorrect and reject the instance. Correctness of these actions
follows if we assume that the partition $\rho$ of $V_s$ we have computed
corresponds to the connected components of $F^\star$, because in the latter case
any solution that connects $w$ to $w'$ will place $w'$ in the same component as
one of $u,v$, and in the former case we are forced to use the selected path, as
otherwise $u,v$ would end up in the same component of $F^\star$.

\item There exist two topo-edges $P_1,P_2$ adjacent to $u_1,v_1\in V_s$ and
$u_2,v_2\in V_s$ respectively, and vertices $w_1\in P_1$ and $w_2\in P_2$ such
that $\{w_1,w_2\}\in D$. If $u_1,v_1,u_2,v_2$ are in four distinct components
of~$F^\star$ we reject the current guess, as it is impossible to place $w_1,w_2$ in
the same component without also placing some of $u_1,v_1,u_2,v_2$ in the same
component.

\item If $u_1,u_2,v_1,v_2,w_1,w_2$ are as previously but $u_1,v_1,u_2,v_2$ are
in three distinct components of $F^\star$, we can assume without loss of
generality that $u_1,u_2$ are in the same component. We replace the demand
$\{w_1,w_2\}$ with the demands $\{w_1,u_1\}$ and $\{w_2,u_2\}$ and reduce to a
previous case.

\end{enumerate}

Finally, if none of the previous cases apply we have arrived at an instance
where all remaining demands $\{w_1,w_2\}\in D$ satisfy that $w_1,w_2$
belong in two distinct topo-edges $P_1,P_2$, which are incident on $u_1,v_1\in
V_s$ and $u_2,v_2\in V_s$ respectively, such that $u_1,u_2$ are in the same
component of $F^\star$, and so are~$v_1,v_2$, but the component of $u_1,u_2$ is
distinct from the component of $v_1,v_2$.  We will find the best way to satisfy
such demands by solving an auxiliary problem.

Fix two sets $C_1,C_2$ of the partition $\rho$ of $V_s$ which we have computed
and consider every topo-edge $P$ with one endpoint in $C_1$ and the other in
$C_2$. We construct a new instance of \SF on a graph $G_2$ by taking the union
of all such topo-edges and then contracting all vertices of $C_1$ into a single
vertex $c_1$ and all vertices of $C_2$ into a single vertex $c_2$. We include
in the new instance all demands with at least one endpoint on one of the
internal vertices of the topo-edges we used; note that such demands also have
the second endpoint in $G_2$. Let $G_1$ be the instance induced from the
original graph if we delete all internal topo-edge vertices which appear in
$G_2$. Note that every demand of the original instance appears in either $G_1$
or $G_2$.

We will now state two claims:

\begin{claim}\label{claim:g1g2} %
If the optimal \SF solution on the
instance $G_2$ constructed above has cost~$b_2$, then we have the following:
$G$ has a solution of cost at most $b$ consistent with the guess $\rho$ if and
only if $G_1$ has a solution consistent with the guess $\rho$ of cost at most
$b-b_2$.  \end{claim}

\begin{claim}\label{claim:g2opt} %
The optimal solution to $G_2$ can be
computed in polynomial time by a reduction to the \textsc{Min Cut} problem.
\end{claim}

Let us explain why the claims are sufficient to conclude our algorithm. We
consider every pair of sets~$C_1,C_2\in\rho$ (of which there are $O(k^2)$) and
for each such pair the claims imply that we can decompose the instance into two
instances $G_1,G_2$, such that $G_2$ can be solved in polynomial time, and
using the optimal value we calculate for $G_2$ we can reduce solving $G$ to
solving $G_1$. Repeating this for all  pairs results in an instance with no
demands. Putting everything together, we have that for one of $2^{5k}$ possible
guesses (on which topo-edges are fully used) we apply a series of
polynomial-time reduction rules that allow us to decompose the instance into
$O(k^2)$ polynomial-time solvable sub-problems. We therefore obtain an exact
algorithm running in $2^{O(k)}n^{O(1)}$ time. \qedhere


%
\begin{proof}[Proof of \cref{claim:g1g2}]

If $G_1$ has a solution of cost $b-b_2$ consistent with $\rho$, then we can form
a solution for $G$ by taking the union of the solution for $G_1$ with an
optimal solution for $G_2$.  This will have cost at most~$b$. Furthermore,
recall that all demands of $G$ appear in either $G_1$ or $G_2$. Demands that
appear in $G_1$ are clearly satisfied by the new solution in $G$, while demands
that appear in $G_2$ are satisfied because the solution in $G_1$ is consistent
with $\rho$, so it contains paths between any two $u,v\in C_1$ for each $C_1\in
\rho$.

For the converse direction, suppose $G$ has a solution of cost $b$ consistent
with $\rho$. We observe that this solution restricted to $G_2$ is a feasible
solution (which furthermore places $c_1,c_2$ in distinct components), hence
must have cost at least $b_2$. Therefore, the solution restricted to edges of
$G_1$ has cost at most $b-b_2$. Because all topo-edges included in $G_2$ are
topo-edges which are not fully used (according to the guess that gave us
$\rho$), the solution we construct in $G_1$ is still consistent with $\rho$ and
satisfies all demands.  \end{proof}
%

%
\begin{proof}[Proof of \cref{claim:g2opt}]

Before we begin, we perform a basic simplification step. If the instance
contains a Steiner vertex $v$ of degree $2$ (that is, a vertex not incident on
any demand), with neighbors $u_1,u_2$, then we remove $v$ from the instance and
add an edge $u_1u_2$ with weight equal to $w(vu_1)+w(vu_2)$. It is not hard to
see that the new instance is equivalent ($v$ would only be used in a solution
if both its incident edges are used), and we now know that all vertices of
degree $2$ are terminals.

Recall that we have a graph $G_2$ with two special vertices $c_1,c_2$ such that
the graph consists of a collection of parallel paths with endpoints $c_1,c_2$,
and furthermore every demand is between two internal vertices of distinct
paths. For the purposes of the larger algorithm we are interested in computing
the best solution where $c_1,c_2$ are in distinct components, but for the sake
of completeness let us briefly note that $G_2$ can be solved to optimality
without this constraint, as the best solution where $c_1,c_2$ are in the same
component is just a minimum cost spanning tree of $G_2$ (here we are using the
fact that all internal vertices of all paths are terminals).

In order to compute the best solution that places $c_1,c_2$ into distinct
components, we will reduce the problem to \textsc{Min Cut}. Let $N$ be a
sufficiently large value, for example set $N$ to be the sum of all edge weights
of the instance. We construct a \textsc{Min Cut} instance on the same graph but
with weight function $w'(e)=N-w(e)$. Furthermore, for all $w_1,w_2$ such that
$\{w_1,w_2\}\in D$ we add an edge $w_1w_2$ and set $w'(w_1w_2)=n^2N$.

Our claim is now that if $F_c$ is a set of edges that gives a minimum weight
$c_1-c_2$ cut in the new instance, then the complement of $F_c$ is a minimum
cost \SF solution for $G_2$ that places $c_1,c_2$ in distinct components.

To prove the claim, suppose that $F_c$ is a minimum-weight $c_1-c_2$ cut in the
new graph. We observe that $F_c$ cannot include any of the edges we added
between the endpoints of demands $(w_1,w_2)\in D$, as such edges have a very
high cost (deleting every other edge would be cheaper). Furthermore, because
all edges have positive weight and the cut $F_c$ is minimal, removing $F_c$
from the graph must leave exactly two connected components, one containing each
of $c_1,c_2$. Hence, for each $(w_1,w_2)\in D$, if we keep in the graph all
edges not in $F_c$, $w_1,w_2$ are in the same component, and we have a feasible
\SF solution for all the demands. In the other direction, consider an optimal
\SF solution that places $c_1,c_2$ in distinct components, and let $F_c'$ be
the set of edges of $G_2$ \emph{not} included in the solution. $F_c'$ must be a
valid $c_1-c_2$ cut, because it contains at least one edge from each
topological edge connecting $c_1$ to $c_2$ (otherwise $c_1,c_2$ would be in the
same component); and as each demand $(w_1,w_2)\in D$ is satisfied, therefore,
$w_1,w_2$ are either in the component of $c_1$ or in the component of $c_2$. We
have therefore established a one-to-one mapping between optimal minimum cuts
and optimal \SF solutions and conclude the claim by observing that by
minimizing the weight of $F_c$ in the \textsc{Min Cut} instance, we are
maximizing the weight of non-selected edges in the \SF instance (thanks to the
modified weight function), hence we are selecting an optimal \SF solution.
\end{proof}
%

\end{proof}

